\chapter{Metricas Fisiologicas y Firmas de Fatiga}
\label{cap:metrics}

\section{Introduction}

Del dataset crudo con multiples archivos a diferentes frecuencias, se extrajeron
\textbf{19 variables fisiologicas agregadas}. Este capitulo explica cada una y,
crucialmente, como se relacionan con fatiga fisica y mental.

\section{Metricas Cardiacas}

\subsection{Frecuencia Cardiaca (HR - bpm)}

\begin{small}
\begin{table}[H]
\centering
\begin{tabular}{|l|p{10cm}|}
\hline
\textbf{Unidad} & Latidos por minuto (bpm) \\
\textbf{Rango normal} & 60-100 bpm (en reposo) \\
\textbf{Rango en dataset} & 58.9-162.3 bpm \\
\textbf{Media dataset} & 88.4 +/- 20.4 bpm \\
\hline
\end{tabular}
\end{table}
\end{small}

\subsubsection{Interpretación}

\begin{itemize}
    \item \textbf{M1 (Baseline)}: HR ~76 bpm - Normal en reposo
    \item \textbf{M2 (Post-ejercicio)}: HR ~105 bpm - AUMENTA con fatiga fisica
    \item \textbf{M3 (Post-mental)}: HR ~84 bpm - Normaliza parcialmente
\end{itemize}

\textbf{Firma de fatiga fisica}: HR aumenta significativamente con intensidad.
Low: 82 bpm, Medium: 95 bpm, High: 108 bpm (tendencia clara).

\subsubsection{Formula}

Simplemente promedio de latidos contados en ventana temporal, derivados de
picos ECG (BioHarness) o PPG (eSense, E4).

\subsection{Variabilidad de Frecuencia Cardiaca (HRV - ms)}

\begin{small}
\begin{table}[H]
\centering
\begin{tabular}{|l|p{10cm}|}
\hline
\textbf{Unidad} & Milisegundos (SDNN - Standard Deviation of NN intervals) \\
\textbf{Rango normal} & 30-100 ms (en reposo, signo de buena salud cardiovascular) \\
\textbf{Rango en dataset} & 11.3-189.4 ms \\
\textbf{Media dataset} & 69.2 +/- 31.2 ms \\
\hline
\end{tabular}
\end{small}
\end{table}

\subsubsection{Interpretación}

HRV mide la variabilidad en espaciamiento entre latidos. Un corazon bien regulado
``sintoniza'' dinámicamente con demandas. HRV baja = inflexibilidad cardiovascular
= estrés/fatiga.

\begin{itemize}
    \item \textbf{HRV Alta (>100 ms)}: Parasimpatico activo = relajacion, recuperacion
    \item \textbf{HRV Moderada (50-100 ms)}: Balance simpático-parasimpático
    \item \textbf{HRV Baja (<30 ms)}: Simpatico dominante = estrés agudo, fatiga
\end{itemize}

\textbf{Firma de fatiga}: HRV DISMINUYE durante M2 post-ejercicio (69->89->49 ms
en M1->M2->M3), indicando respuesta estresante.

\subsubsection{Formula}

$$\text{HRV (SDNN)} = \sqrt{\frac{1}{N-1} \sum_{i=1}^{N} (NN_i - \overline{NN})^2}$$

Donde $NN_i$ es la duracion del i-esimo intervalo entre latidos sucesivos (R-R).

\section{Metricas Respiratorias}

\subsection{Tasa Respiratoria (BR - brpm)}

\begin{small}
\begin{table}[H]
\centering
\begin{tabular}{|l|p{10cm}|}
\hline
\textbf{Unidad} & Respiraciones por minuto (brpm) \\
\textbf{Rango normal} & 12-20 brpm (en reposo) \\
\textbf{Rango dataset} & 11.7-33.4 brpm \\
\textbf{Media dataset} & 20.9 +/- 4.6 brpm \\
\hline
\end{tabular}
\end{small}
\end{table}

\subsubsection{Interpretación}

\begin{itemize}
    \item M1 (Baseline): BR ~16.5 brpm - Normal
    \item M2 (Post-ejercicio): BR ~18.9 brpm - Elevada con ejercicio
    \item M3 (Post-mental): BR ~15.5 brpm - Normalizacion
\end{itemize}

Respiracion rapida puede indicar estrés o hiperventilacion (fatiga mental).
Respiracion lenta (si anormal) puede indicar somnolencia severa.

\section{Metricas Electrodermal}

\subsection{Actividad Electrodermal (EDA - µS)}

\begin{small}
\begin{table}[H]
\centering
\begin{tabular}{|l|p{10cm}|}
\hline
\textbf{Unidad} & Microsiemens (conductancia de piel) \\
\textbf{Rango normal} & 0.5-2 µS (basalmente) \\
\textbf{Rango dataset} & 0.08-7.68 µS \\
\textbf{Media dataset} & 1.30 +/- 1.82 µS \\
\hline
\end{tabular}
\end{small}
\end{table}

\subsubsection{Interpretación}

EDA refleja activacion del sistema nervioso simpatico. Medida muy sensible a
emociones, estrés y arousale cognitivo.

\begin{itemize}
    \item \textbf{EDA Baja}: Desapego mental, somnolencia
    \item \textbf{EDA Moderada}: Concentracion normal
    \item \textbf{EDA Alta}: Estrés, atencion, respuesta emocional
\end{itemize}

En FatigueSet: M1 (0.54 µS) -> M2 (1.68 µS) -> M3 (1.66 µS).
Aumenta en M2 reflejando demanda de ejercicio.

\section{Metricas Cerebrales: Bandas EEG}

\subsection{EEG Alpha (8-12 Hz)}

\textbf{Rango}: -0.5 a +1 Bels (escala logaritmica).
\textbf{Interpretacion}: Mayor en relajacion, menor en concentracion.

En dataset: Media 0.575 +/- 0.159 Bels.

\subsubsection{Firma de Fatiga}
Alpha aumentado puede indicar somnolencia si es excesivo.
En FatigueSet: Disminuye ligeramente en M2 (actividad), vuelve en M3.

\subsection{EEG Beta (12-30 Hz)}

\textbf{Rango}: Similar a Alpha.
\textbf{Interpretacion}: Mayor en atención activa, menor en relajacion.

En dataset: Media 0.402 +/- 0.144 Bels.

\subsubsection{Firma de Fatiga}
Beta bajo es sugestivo de bajo arousale o fatiga cognitiva.

\subsection{EEG Theta (4-8 Hz)}

\textbf{Rango}: Similar.
\textbf{Interpretacion}: Elevado en somnolencia, meditacion profunda.

En dataset: Media 0.436 +/- 0.198 Bels.

\subsubsection{Firma de Fatiga}
\textbf{THETA AUMENTADO es el marcador EEG mas fuerte de fatiga mental y somnolencia}.

Patrones esperados:
\begin{itemize}
    \item Fatiga mental severa: Theta muy elevado (>0.6 Bels)
    \item Alerta normal: Theta moderado (0.3-0.4 Bels)
\end{itemize}

\section{Tabla Integrativa: Interpretación por Fase}

\begin{table}[H]
\centering
\caption{Cambios esperados de metricas entre fases}
\begin{tabular}{|l|l|l|l|}
\hline
\textbf{Metrica} & \textbf{M1 (Baseline)} & \textbf{M2 (Post-Ejerc.)} & \textbf{M3 (Post-Mental)} \\
\hline
HR & Baja (~76) & ALTA (~105) & Moderada (~84) \\
HRV & Moderada (~70) & BAJA (~89) & Baja (~49) \\
BR & Normal (~16) & ELEVADA (~19) & Normal (~15) \\
EDA & Baja (~0.54) & ELEVADA (~1.68) & Elevada (~1.66) \\
EEG Alpha & Moderada & Ligeramente baja & Aumenta \\
EEG Beta & Moderada & Puede aumentar & Aumenta \\
EEG Theta & Baja & Ligeramente elevada & Elevada (cansancio) \\
\hline
\end{tabular}
\end{table}

\section{Resumen: 19 Variables Agregadas}

Las 19 columnas del archivo fatigueset\_aggregated\_features.csv son:

\begin{enumerate}
    \item participante, sesion, intensidad, intensidad\_num, fase, fase\_num (6 identificadores)
    \item fatiga\_fisica, fatiga\_mental (2 autoreportes)
    \item hr\_media, hr\_std (2 metricas cardiaca)
    \item br\_media, br\_std (2 metricas respiratoria)
    \item hrv\_media, hrv\_std (2 metricas variabilidad cardiaca)
    \item eda\_media, eda\_std (2 metricas electrodermal)
    \item eeg\_alpha\_media, eeg\_beta\_media, eeg\_theta\_media (3 metricas cerebrales)
\end{enumerate}

Total: 6 + 2 + 2 + 2 + 2 + 2 + 3 = 19 columnas en el archivo agregado.

Cada metrica (excepto autoreportes categóricos) esta disponible como media y
desviacion estandar, permitiendo evaluar estabilidad de la senal.
